\section{Análisis de caso: de la descripción de la tarea a la trayectoria ejecutable}

\subsection{El ejemplo de ``reservar un restaurante que cumpla las restricciones''}
La clase recurre varias veces a tareas de búsqueda de restaurantes para ilustrar las dificultades de ejecución del agente. Este tipo de tareas parece simple, pero en páginas web reales suele incluir restricciones múltiples: ubicación, precio, calificación, horario de atención, preferencia de cocina y el proceso de pedido. Si el sistema no cuenta con una gestión explícita del estado, es fácil que olvide condiciones a la mitad del proceso.

\begin{importantbox}{Estructura recomendada para descomponer la tarea}
\begin{enumerate}
    \item \textbf{Extracción de restricciones}: convertir los requisitos en lenguaje natural en campos estructurados;
    \item \textbf{Recuperación de candidatos}: generar un conjunto de candidatos dentro del sitio o entre sitios;
    \item \textbf{Verificación de evidencia}: verificar una por una las restricciones duras (como el tope de precio);
    \item \textbf{Ejecución de la acción}: entrar al proceso de pedido y completar el llenado de los campos clave;
    \item \textbf{Confirmación final}: confirmar por partida doble, mediante el estado de la página y el verifier, que la tarea se completó.
\end{enumerate}
\end{importantbox}

\begin{knowledgebox}{Por qué son clave las ``restricciones estructuradas''}
Si las restricciones no se estructuran, el agente tiende a perder información clave a lo largo de una cadena larga y solo conserva preferencias vagas. Al almacenar las restricciones explícitamente como pares clave-valor, puede hacerse una verificación de consistencia antes de cada paso de ejecución, lo que reduce de forma notable el costo de tener que rehacer trabajo por ``descubrir hasta el final que no se cumplía la condición''.
\end{knowledgebox}

\subsection{Control de calidad a nivel de trayectoria}
En esta tarea, que un paso individual sea correcto no significa que la trayectoria sea correcta. Por ejemplo, el modelo puede elegir primero un restaurante que se ve bien, pero cuyo precio u horario disponible no cumple los requisitos. El control de calidad a nivel de trayectoria debe verificar la ``consistencia con el objetivo'', y no solo la ``ejecutabilidad de la acción''.

\begin{table}[H]
\centering
\small
\begin{tabular}{p{2.8cm}p{3.5cm}p{4.4cm}}
\toprule
\textbf{Capa de verificación} & \textbf{Objeto de verificación} & \textbf{Errores comunes} \\
\midrule
Consistencia de campos & Precio, horario, ubicación, calificación & Campos omitidos, confusión de unidades, restricciones en conflicto \\
Consistencia de página & Si la página actual contiene la evidencia buscada & Error al identificar el elemento, lectura incorrecta de la información de la tarjeta \\
Consistencia del proceso & Si se entró a la ruta correcta de pedido & Sitio equivocado, entrada equivocada, detención prematura \\
Consistencia final & Si el resultado de la tarea coincide con el requisito original & Éxito parcial pero incumplimiento global \\
\bottomrule
\end{tabular}
\caption{Control de calidad a nivel de trayectoria en la tarea de búsqueda de restaurantes}
\end{table}

\begin{figure}[H]
\centering
\includegraphics[width=0.82\textwidth]{frame-07-search-eval.jpg}
\caption{Video aproximadamente en 00:34:39: comparar entre trayectorias candidatas antes de seguir expandiendo, lo que refleja la estrategia de ``evaluar antes de ejecutar''}
\end{figure}

\subsection{Plantillas de ejecución reutilizables}
Para reducir la necesidad de explorar desde cero en cada tarea, puede mantenerse una ``biblioteca de plantillas de tareas''. La plantilla no es un guion fijo, sino un ``esqueleto de proceso que puede restringirse'': define el orden de las fases, los puntos de verificación clave y la estrategia de reversión. Así se conserva la flexibilidad y, a la vez, se reduce de forma notable la varianza en tareas largas.

\begin{warningbox}{Los límites de usar plantillas}
Una plantilla demasiado rígida perjudica la generalización, y una demasiado laxa no logra restringir los errores. Lo razonable es usar la plantilla en los pasos de alto riesgo (pago, identidad, página de confirmación) y dejar la exploración para los pasos de bajo riesgo (recuperación y ordenamiento de candidatos).
\end{warningbox}

\subsection{Resumen de la sección}
Lo esencial en las tareas web complejas no es ``lograr que el modelo adivine mejor'', sino ``lograr que la trayectoria sea más controlable''. Las restricciones estructuradas, la verificación a nivel de trayectoria y las plantillas reutilizables son los tres trabajos básicos para convertir un sistema de demostración en un sistema estable.
